\chapter*{Kurzfassung}

Diese Arbeit entwickelt datengetriebene Indikatoren für die vorausschauende Instandhaltung
des Pneumatiksystems eines Metro-Zuges auf Basis des MetroAT-Datensatzes: ein Jahr an
Sensoraufzeichnungen mit 1\,Hz und Min-Max-Normalisierung von einem einzelnen Zug der
Wiener Linien. Die Analyse folgt einem methodischen Grundprinzip: der Vermeidung von
Leakage bzw. Zirkularität. Betriebszustände werden ausschließlich aus der Bewegung
definiert, sodass die Sensoren des Luftsystems, die nicht in diese Definition eingehen, als
unabhängige Evidenz erhalten bleiben.

Das Clustern der Bewegungsmerkmale liefert vier Betriebszustände: Stillstand (48\%),
Bremsen (20\%), Fahren (17\%) und Beschleunigen (14\%). Die Clusteranzahl $k=4$ wurde über
die größte Verbesserung des Informationskriteriums gewählt und erwies sich unter
Resampling als stabil. Die Beschränkung der Zustandsdefinition auf die Bewegung korrigierte
einen früheren Fehler, bei dem ein stehender Zug mit angelegter Bremse fälschlich als
Bremsen eingestuft wurde; der korrigierte Bremszustand ist zu 99,9\% echt. Auf
zurückgehaltenen Daten erkennen die Luftsystem-Sensoren den Bremszustand mit einer ROC-AUC
von 0,95, und von 0,92 selbst dann, wenn das Bremssignal weggelassen und nur unabhängige
Zustandssensoren verwendet werden. Das Bremsen hinterlässt also einen Fingerabdruck im
gesamten Pneumatiksystem.

Die Bremsintensität wurde auf Klassenstruktur geprüft und erwies sich als Kontinuum, nicht
als Menge diskreter Klassen. Ein Gradient-Boosted-Tree erklärt außerhalb der Stichprobe
einen mäßigen Anteil der Intensität aus den Nicht-Geschwindigkeitssensoren
($R^2 \approx 0{,}43$--$0{,}47$, gegen eine Mittelwert-Baseline) und rund 83\% der pro Halt
dissipierten spezifischen Energie. Schließlich wurde das Bremsverhalten auf ein Signal vor
dokumentierten Bremssystem-Ausfällen untersucht. Über Effektstärken-, Change-Point- und
Energieanalysen hinweg ist das Ergebnis nahezu null; mit nur acht Bremsausfällen ist die
Studie explorativ und liefert keinen einsetzbaren Ausfalldetektor. Das negative Ergebnis
legt nahe, dass sich Bremsfehler eher während der Lade- und Ruhephasen als während des
Bremsens zeigen, also in die Richtung, die eine Folgeuntersuchung einschlagen sollte.
